\documentclass[11pt]{article}
\usepackage[T1]{fontenc}
\usepackage{lmodern}
\usepackage[margin=1in]{geometry}
\usepackage{amsmath,amssymb,amsthm,booktabs}
\usepackage[colorlinks=true,urlcolor=blue,linkcolor=blue,citecolor=blue]{hyperref}
\usepackage{microtype}
\hypersetup{pdftitle={A derivative determinant bound for the tenth discriminant},
 pdfauthor={Research note prepared with Codex},
 pdfsubject={Version 0.4.0: q < 4000}}
\newtheorem{theorem}{Theorem}
\newtheorem{lemma}[theorem]{Lemma}
\newtheorem{proposition}[theorem]{Proposition}
\newcommand{\Z}{\mathbb Z}
\newcommand{\Q}{\mathbb Q}
\newcommand{\C}{\mathbb C}
\newcommand{\OK}{\mathcal O_K}
\newcommand{\Tr}{\operatorname{Tr}}
\newcommand{\Res}{\operatorname{Res}}
\newcommand{\ord}{\operatorname{ord}}
\title{A derivative determinant bound for the tenth discriminant\\
\large Version 4: \(q<4\,000\)}
\author{Research note prepared with Codex}
\date{Version 0.4.0 --- 20 October 2026}
\begin{document}
\maketitle

\begin{abstract}
We obtain the analytic cutoff \(q<4\,000\) for imaginary quadratic fields
of class number one, a factor of 250 below the preceding cutoff \(10^6\).
The direct split-prime estimate of v3 is combined with Green's derivative
determinant, replacing Banks's quartic-field construction. We sharpen the
nonsingularity criterion, choose points in a rectangle adapted to the
quadratic integer lattice, and keep the decrease of polynomial degree
under differentiation in the absolute-value estimate. The determinant
argument is proved here, including the resultant estimate. Classical
explicit prime estimates remain external inputs. An exact rational
certificate checks the constants and uniform contradiction; an
exhaustive integer check below the cutoff leaves the nine familiar
discriminants. No earlier version is needed to read the proof.
\end{abstract}

\paragraph{Status and attribution.}
This is an AI-assisted research draft.
The new v4 mathematical argument has not yet had an independent
cross-check. The supplied programs check arithmetic and finite cases;
they do not formally verify the proof. The derivative determinant and
the resultant method originate in Green's note \cite{green}. The
refinements and explicit parameter choices below were developed with
Codex. Both v3 papers have passed independent mathematical checks;
those checks do not cover the new argument in this version.

\section{Result, notation, and proof strategy}

Let \(D<0\) be a fundamental discriminant, let
\(K=\Q(\sqrt D)\), and put \(q=|D|\). The class number \(h(K)\)
is the number of ideal classes of its ring of integers \(\OK\).
Thus \(h(K)=1\) means that every nonzero integral ideal is principal.
All logarithms in this paper are natural.

\begin{theorem}\label{main}
If \(h(K)=1\), then \(q<4\,000\).
The finite verification supplied with this paper consequently gives
\[
 h(K)=1\quad\Longleftrightarrow\quad
 q\in\{3,4,7,8,11,19,43,67,163\}.
\]
\end{theorem}

The analytic cutoff, rather than the known classification by other
methods, is the object of Tao's challenge \cite{tao}.
The argument has four steps.
\begin{enumerate}
\item Under class number one, split primes are norms of primitive
quadratic integers. Counting these integers bounds the weighted sum
of split primes, so an explicit total prime sum forces many inert primes.
\item We evaluate monomials and their derivatives at a finite set of
quadratic integers. A resultant argument proves that the associated
integer Gram determinant is positive.
\item Frobenius at every inert prime forces a large prime-power divisor
of this determinant. Its absolute value has a competing upper bound.
For every \(q\ge4\,000\), the two bounds are inconsistent.
\item A separate, small integer computation excludes the remaining
candidates except the nine stated in the theorem.
\end{enumerate}

The proof uses quadratic ideal theory, elementary polynomial algebra,
and the two Rosser--Schoenfeld estimates stated in
Section~\ref{primes}. Reduction of positive binary quadratic forms is
used only in the finite classification. The new determinant argument is
included in full. We use neither a Heegner method nor computed
zero-free regions for quadratic Dirichlet \(L\)-functions.

To prove the cutoff, assume from now until Section~\ref{finite} that
\[
 h(K)=1,\qquad q\ge4\,000,\qquad L=\log q,\qquad u=\sqrt q.
\]
The symbols \(R,S,J\) below denote determinant parameters;
in particular, \(R\) does \emph{not} denote \(\sqrt q\).

\section{Class number one and a direct bound for split primes}

For a prime \(p\), the quadratic character \(\chi_D(p)\) is respectively
\(1,-1,0\) when \(p\) splits, is inert, or ramifies in \(K\).
In the split or ramified case there is an integral ideal of norm \(p\).
Under \(h(K)=1\), a generator \(\alpha\) of that ideal satisfies
\(|N_{K/\Q}(\alpha)|=p\).

In the standard integral basis \(1,\omega\), one has
\(\operatorname{Im}\omega=u/2\) and
\(\operatorname{Re}\omega\in\{0,1/2\}\). Hence every nonrational
quadratic integer has norm at least \(q/4\); a rational integer has a
square norm. It follows that
\begin{equation}\label{small-inert}
 p<q/4\quad\Longrightarrow\quad\chi_D(p)=-1.
\end{equation}
In particular, 2 is inert. The quadratic character at 2 then implies
\(D\equiv5\pmod8\), or \(q\equiv3\pmod8\). We may therefore use
\[
 \omega=\frac{1+i\sqrt q}{2},\qquad
 \OK=\Z+\Z\omega,\qquad
 Q(a,b)=|a+b\omega|^2=(a+b/2)^2+qb^2/4.
\]

\subsection{The primitive lattice count}

Let \(A(x)\) count pairs \((a,b)\in\Z\times\Z_{\ge1}\) with
\(Q(a,b)\le x\). For fixed \(b\), the allowable real interval for \(a\)
has length
\(f(b)=2\sqrt{x-qb^2/4}\), when this is defined. Its integer-point
count differs from its length by at most one. Since \(f\) is decreasing,
comparison with its integral gives
\begin{equation}\label{raw-count}
 \frac{\pi x}{u}-2\sqrt x-\frac{2\sqrt x}{u}
 \ \le A(x)\le\frac{\pi x}{u}+\frac{2\sqrt x}{u}.
\end{equation}
Here \(\int_0^{2\sqrt x/u}f(t)\,dt=\pi x/u\), and the lower
integral comparison loses at most \(f(0)=2\sqrt x\).

Sieve out common factors of \(a,b\) among \(2,3,5,7\). With \(P=210\),
the resulting count is
\[
 A_P(x)=\sum_{d\mid P}\mu(d)A(x/d^2).
\]
Inserting the appropriate upper or lower estimate from
\eqref{raw-count}, according to the sign of \(\mu(d)\), gives
\begin{equation}\label{sieved-count}
 A_P(x)\le\frac{\pi\rho x}{u}+B_q\sqrt x,
 \qquad
 \rho=\frac{768}{1225},\qquad
 B_q=\frac{88}{35}+\frac{192}{35u}.
\end{equation}
Indeed, the three elementary products used here are
\[
 \prod_{p\mid P}(1-p^{-2})=\frac{768}{1225},\quad
 \prod_{p\mid P}(1+p^{-1})=\frac{96}{35},\quad
 \prod_{p\mid P}(1-p^{-1})=\frac{8}{35}.
\]
The coefficient of \(\sqrt x\) is the second product minus the third,
plus \(2/u\) times the second.

A split prime \(p\) has two distinct prime ideals of norm \(p\).
Their generators give exactly two points with norm \(p\) and \(b>0\):
the only units for \(q>4\) are \(\pm1\), and conjugation exchanges the
two prime ideals. Each such point has \(\gcd(a,b)=1\), since otherwise
a nontrivial square would divide \(p\). Thus, if \(\pi_+(x)\) counts
split primes up to \(x\), then \(2\pi_+(x)\le A_P(x)\).

\subsection{Partial summation}

Set
\[
 S_+(X)=\sum_{\substack{p\le X\\\chi_D(p)=1}}\frac{\log p}{p},
 \qquad a_0=q/4.
\]
There are no split primes below \(a_0\). Since \(a_0>e\), the function
\(\log t/t\) is decreasing throughout the interval of integration.
Partial summation applied to \eqref{sieved-count} gives, for \(X\ge a_0\),
\begin{equation}\label{split-bound}
 S_+(X)\le
 \frac{\pi\rho}{4u}
 \bigl((\log X)^2-(\log a_0)^2+2\log a_0\bigr)
 +\frac{2B_q(\log a_0+1)}{u}.
\end{equation}
For completeness, the linear count contributes the first term exactly.
The square-root count contributes
\[
 \frac{B_q(\log a_0+1)}{\sqrt{a_0}}
 -\frac{B_q(\log X+2)}{2\sqrt X};
\]
discarding its last, negative term yields \eqref{split-bound}.

Every ramified prime divides \(q\), and \eqref{small-inert} also
applies to ramified primes. Therefore
\begin{equation}\label{ramified}
 S_0(X):=\sum_{\substack{p\le X\\p\mid q}}\frac{\log p}{p}
 \le\frac4q\sum_{p\mid q}\log p\le\frac{4L}{q}.
\end{equation}

\section{The explicit ordinary-prime estimates}\label{primes}

Write
\(A_{\mathrm{pr}}(x)=\sum_{p\le x}\log p/p\) and
\(\vartheta(x)=\sum_{p\le x}\log p\).
Theorems 6 and 9 of Rosser--Schoenfeld \cite{rs} give a real constant
\(c_0\) such that
\begin{equation}\label{rs}
 |A_{\mathrm{pr}}(x)-\log x-c_0|<\frac1{2\log x}\quad(x\ge319),
 \qquad
 \vartheta(x)<cx\quad(x>0),\quad c=1.01624.
\end{equation}
Evaluation of the finite prime sum at 319, with rigorous logarithm
intervals, gives
\[
 \log319-A_{\mathrm{pr}}(319)+\frac1{2\log319}<1.334.
\]
Consequently \(c_0>-1.334\), so
\begin{equation}\label{prime-total}
 A_{\mathrm{pr}}(x)>\log x-1.334-\frac1{2\log x}\qquad(x\ge319).
\end{equation}
All finite numerical inequalities in this paper are checked by the
accompanying exact arithmetic certificate.

We will separate the primes at \(H=50\). Define
\begin{align*}
 a_H&=\sum_{p\le50}\frac{\log p}{p},&
 \theta_H&=\sum_{p\le50}\log p,\\
 W&=\sum_{p\le50}\left(\frac1{50}-\frac{p}{7500}\right)\log p.
\end{align*}
The required bounds are
\begin{equation}\label{small-constants}
 a_H<2.77035,\qquad40.9602<\theta_H<41,\qquad W>0.6732.
\end{equation}
Integration by parts in \eqref{rs} also gives
\begin{equation}\label{prime-tail}
 \sum_{p>50}\frac{\log p}{p^2}
 =-\frac{\theta_H}{50^2}
   +2\int_{50}^\infty\frac{\vartheta(t)}{t^3}\,dt
 < b:=\frac{2c}{50}-\frac{40.9602}{50^2}=0.02426552.
\end{equation}
Finally, another partial summation gives
\(A_{\mathrm{pr}}(x)\le c(1+\log(x/2))<2\log x\) for \(x\ge2\).

\section{A sharpened derivative interpolation lemma}\label{rank}

This section adapts Green's resultant proof \cite{green}. We retain the
actual constant in the bound for the number of points on a curve, and
retain the loss of \(w\)-degree when removing factors depending only on
\(w\). These changes give the following useful criterion.

\begin{lemma}\label{rank-lemma}
Let \(P_i=(z_i,w_i)\in\C^2\), \(1\le i\le n\), with all \(w_i\)
distinct. Suppose every nonzero polynomial \(F\in\C[z,w]\) satisfies
\[
 \#\{i:F(P_i)=0\}\le A\deg F.
\]
Let \(R,S,J\) be positive integers with
\begin{equation}\label{rank-condition}
 J\le2R+A,\qquad nJ>2RS+A(R+S).
\end{equation}
Then the \(RS\) vectors
\[
 v_{r,s}=\bigl(\partial_z^j(z^r w^s)(P_i)\bigr)_{
                 1\le i\le n,\ 0\le j<J},
 \qquad0\le r<R,\quad0\le s<S,
\]
are linearly independent over \(\C\).
\end{lemma}

\begin{proof}
A dependence would give a nonzero polynomial of degrees less than
\((R,S)\) whose first \(J\) derivatives in \(z\) vanish at every point.
Factor it over \(\C\) as
\[
 F(z,w)=F_0(w)\prod_{\ell=1}^k F_\ell(z,w),
\]
where each \(F_\ell\) is irreducible and has positive \(z\)-degree
\(d_\ell\); factors may be repeated. Write \(d'_\ell=\deg_w F_\ell\)
and \(d'_0=\deg F_0\). We can use the weak bounds
\(\sum d_\ell\le R\) and \(\sum d'_\ell\le S-d'_0\).

For an irreducible factor of bidegree \((d,d')\) with \(d>0\), the
polynomial \(T(w)=\Res_z(F_\ell,\partial_zF_\ell)\) is nonzero and
has degree at most \(2dd'\). No specialization
\(F_\ell(z,w_i)\) is identically zero: otherwise \(w-w_i\) would
divide the irreducible polynomial. At a root of multiplicity \(m\),
\begin{equation}\label{resultant-order}
 m\le1+\ord_{w=w_i}T(w).
\end{equation}
Here is a direct justification. The Sylvester matrix has size \(2d-1\).
After specialization at \(w_i\), both polynomials are divisible by
\((z-z_i)^{m-1}\), so the matrix has rank deficiency at least \(m-1\).
A constant invertible row transformation can make that many rows zero
at \(w_i\). Each such row is then divisible by \(w-w_i\), so the
determinant vanishes to order at least \(m-1\). Its entries have
\(w\)-degree at most \(d'\), which also proves the degree bound.
Nonvanishing follows from irreducibility and characteristic zero.

The number of points where this factor vanishes is at most
\(A(d+d')\). Summing \eqref{resultant-order} over these points, using
distinct \(w_i\), bounds the sum of its specialized multiplicities by
\(A(d+d')+2dd'\).

There are at most \(d'_0\) indices where \(F_0(w_i)=0\). At every
other index the specialization of \(F\) has multiplicity at least \(J\).
Thus
\begin{align*}
 J(n-d'_0)
 &\le A(R+S-d'_0)+2R(S-d'_0),\\
 nJ
 &\le A(R+S)+2RS+(J-A-2R)d'_0
 \le A(R+S)+2RS.
\end{align*}
The last inequality uses \(J\le2R+A\), and contradicts
\eqref{rank-condition}.
\end{proof}

\section{The integral determinant and its prime divisibility}\label{det}

Take distinct points \(z_i\in\OK\) and \(w_i=\overline{z_i}\) satisfying
Lemma~\ref{rank-lemma}. The derivative vectors lie in \(\OK^{nJ}\).
Form the real Gram matrix
\[
 G_{(r,s),(r',s')}=
 \Tr_{K/\Q}\!\left(\sum_{i,j}
 (v_{r,s})_{i,j}\,\overline{(v_{r',s'})_{i,j}}\right).
\]
Every entry is an integer. The trace pairing on this imaginary
quadratic field is the positive real inner product
\(2\operatorname{Re}(\alpha\overline\beta)\). Complex linear
independence of the vectors implies their real linear independence.
Consequently
\begin{equation}\label{positive-det}
 \Delta:=\det G\in\Z_{>0}.
\end{equation}

\subsection{The absolute-value bound}

Suppose \(|z_i|=|w_i|\le T\) and \(1\le T\le R\). An entry involving
the \(j\)-th derivative has absolute value at most
\[
 R^j T^{r+s-j}\le T^{r+s}(R/T)^{J-1};
\]
when \(j>r\), it is zero. Hadamard's inequality for the Gram matrix
therefore yields
\begin{equation}\label{arch}
 \frac{\log\Delta}{R^2S}
 \le\frac{\log(2nJ)}R
 +\left(1+\frac SR\right)\log T
 +\frac{2J}R\log(R/T).
\end{equation}
In obtaining this slightly weaker expression, we used
\(\sum_{r,s}2(r+s)=RS(R+S-2)\) and then discarded negative terms.
Retaining the factor \(T^{-j}\) from differentiation is useful here.

\subsection{The inert-prime bound, including the prime 2}

If \(p\) is inert, \(\OK/p\OK\) is the field of \(p^2\) elements.
Its Frobenius is the nontrivial quadratic automorphism, so
\(z_i^p\equiv\overline{z_i}=w_i\pmod{p\OK}\).
This applies to \(p=2\) as well as odd inert primes because we work
in the full ring of integers. Put \(\psi=z^p-w\).
Its value and all its positive \(z\)-derivatives at these points are
in \(p\OK\): every nonzero derivative coefficient contains \(p\).

For each pair \((r,s)\), let
\(h_{r,s}=\min(\lfloor r/p\rfloor,S-1-s)\).
Replace its polynomial by
\(z^{r-ph_{r,s}}\psi^{h_{r,s}}w^s\).
The binomial expansion has leading term \(z^r w^s\) and other terms
\(z^{r-pi}w^{s+i}\), all with smaller \(r\) and still within the
rectangle of indices. The induced row transformation is an integer
triangular matrix with determinant one. Every derivative of the
replacement polynomial is in \(p^{h_{r,s}}\OK\), by the product rule.
Applying the same transformation to the Gram matrix and factoring
from rows and columns proves
\begin{equation}\label{valuation}
 v_p(\Delta)\ge
 2\sum_{0\le r<R,\ 0\le s<S}\min(\lfloor r/p\rfloor,S-1-s).
\end{equation}
No division by derivative factorials is made in this argument.

\subsection{Explicit lower bounds for the valuation sum}

Let \(e(R,S,p)\) denote the right side of \eqref{valuation}, divided
by \(R^2S\). Counting the pairs for which the minimum is at least
\(h\) gives the exact identity
\begin{equation}\label{layers}
 e(R,S,p)=\frac2{R^2S}
 \sum_{h=1}^{S-1}(S-h)\max(R-ph,0).
\end{equation}
When \(p(S-1)\le R-1\), summing this quadratic expression gives
\begin{equation}\label{small-e}
 e(R,S,p)=\frac{S-1}{R}-\frac{p(S^2-1)}{3R^2}.
\end{equation}
For a bound valid for all \(p\), use
\(\sum_{s=0}^{S-1}\min(k,S-1-s)\ge Sk-k(k+1)/2\).
Writing \(R=ap+b\), \(0\le b<p\), shows
\[
 \sum_{r=0}^{R-1}\lfloor r/p\rfloor
 =\frac{R(R-p)}{2p}+\frac{b(p-b)}{2p}
 \ge\frac{R(R-p)}{2p}.
\]
Also \(\lfloor r/p\rfloor(\lfloor r/p\rfloor+1)
\le r^2/p^2+r/p\). Substitution gives
\begin{equation}\label{tail-e}
 e(R,S,p)\ge
 \left(1-\frac{R-1}{2RS}\right)\frac1p
 -\frac1R-\frac{(R-1)(2R-1)}{6RS\,p^2}.
\end{equation}

\section{Parameters and an adapted lattice rectangle}\label{params}

Choose
\begin{equation}\label{parameters}
 R=q^4,\qquad S=\lfloor R/50\rfloor,\qquad
 J=\lfloor R/55\rfloor.
\end{equation}
Let
\[
 X_0=\frac{3}{2\sqrt2}q^{9/4},\qquad
 Y_0=\frac{3}{\sqrt2}q^{7/4},\qquad
 N_x=\lceil X_0\rceil,\quad N_y=\lceil Y_0\rceil.
\]
Choose \(N_y\) consecutive integers \(b\), centered as closely as
possible about zero, so \(|b|\le N_y/2\). For each such \(b\), choose
\(N_x\) consecutive integers \(a\) so that
\(|a+b/2|\le N_x/2\). Such a choice exists for either parity of \(b\).
Use the \(n=N_xN_y\) points \(z=a+b\omega\).
They satisfy
\begin{equation}\label{point-size}
 |z|^2\le\frac{N_x^2}{4}+\frac{qN_y^2}{16}.
\end{equation}

Both unrounded side lengths exceed \(q^{3/2}\ge240\,000\). Thus
\(N_x<1.0001X_0\), \(N_y<1.0001Y_0\), and
\begin{equation}\label{grid-bounds}
 \frac94R\le n<3R,\qquad
 A:=N_x\le2q^{9/4},\qquad
 |z|\le T:=1.0001\,\frac34q^{9/4}<R.
\end{equation}
Here \(N_x\ge N_y\), since \(X_0/Y_0=\sqrt q/2>1\).

We check the zero-count hypothesis in Lemma~\ref{rank-lemma}, including
the fact that the horizontal integer intervals may depend on \(b\).
The invertible linear substitution
\((a,b)\mapsto(a+b\omega,a+b\overline\omega)\) preserves total degree.
Suppose a polynomial of total degree \(d\) vanishes identically on
\(k\) of the chosen horizontal lines. Factoring out these \(k\)
linear factors leaves a polynomial of degree at most \(d-k\).
On each other line it has at most \(d-k\) zeros. The total is at most
\[
 kN_x+(N_y-k)(d-k)\le N_xd=Ad.
\]
The coordinates \(w_i=\overline{z_i}\) are distinct because the
quadratic integer representation is unique.

The first rank condition, \(J\le2R+A\), is immediate.
For the second, use \(n\ge9R/4\), \(J\ge R/55-1\),
\(S\le R/50\), and \(A\le2q^{9/4}\). It suffices that
\[
 \frac1{1100}>
 \frac{9}{4q^4}+\frac{51}{25q^{7/4}}.
\]
This holds uniformly for \(q\ge4\,000\), even on replacing
\(q^{7/4}\) by the smaller bound \(240\,000\).
Thus \(\Delta\) is a positive integer.

\section{The uniform contradiction}\label{contradiction}

\subsection{Upper bound for the determinant}

Set
\[
 \beta=1+\frac1{50}-\frac2{55}=\frac{541}{550}.
\]
In \eqref{arch}, increasing \(S/R\) to \(1/50\) and \(J/R\) to
\(1/55\) increases the right side, since \(T\ge1\) and \(R/T>1\).
Using \eqref{grid-bounds} and \(\log(1.0001)<0.0001\), we get
\begin{equation}\label{arch-final}
 \frac{\log\Delta}{R^2S}
 <\frac{5189}{2200}L+\beta\log(3/4)+0.000101.
\end{equation}
Indeed \(\log(2nJ)/R\le(2+8L)/q^4<10^{-6}\).

\subsection{Lower bound for the determinant}

Write
\[
 I=\sum_{\substack{50<p\le R\\\chi_D(p)=-1}}\frac{\log p}{p}.
\]
All primes up to 50 are inert by \eqref{small-inert}.
For these primes \eqref{small-e} applies, because their largest value
is 47 and \(S\le R/50\). Since \(S\ge R/50-1\), their total
contribution to \(\log\Delta/(R^2S)\) is at least
\[
 W-\frac{2\theta_H}{R}>W-10^{-6}.
\]
For the larger inert primes, \eqref{tail-e} gives the lower bound
\[
 I-\frac{I}{2S}-\frac{\vartheta(R)}R
       -\frac{R}{3S}\sum_{p>50}\frac{\log p}{p^2}.
\]
We have \(R/S\le50.001\), and
\(I/(2S)<\log R/S\le50.001(4L)/q^4<10^{-6}\).
Combining the bounds, using \eqref{positive-det}, \eqref{valuation},
\eqref{rs}, and \eqref{prime-tail}, yields
\begin{equation}\label{padic-final}
 \frac{\log\Delta}{R^2S}
 >I+W-c-16.667b-0.000002.
\end{equation}

By separating the three splitting types and using \eqref{prime-total}
and \eqref{ramified},
\[
 I>4L-1.334-\frac1{8L}-a_H-S_+(q^4)-\frac{4L}{q}.
\]
Comparison with \eqref{arch-final} forces
\begin{equation}\label{necessary}
 \delta:=\frac{3611}{2200}
 <\frac{S_+(q^4)}L+\frac4q+\frac KL+\frac1{8L^2},
\end{equation}
where we may use the following rational upper bound for the constant:
\begin{align}
 K={}&1.334+2.77035+c+16.667b-0.6732
       -\beta(0.28768)+0.001\label{k}\\
 ={}&\frac{628354500503}{137500000000}<4.570.\nonumber
\end{align}
Here \(\log(4/3)>0.28768\), and \(0.001\) exceeds all the losses
\(0.000101+0.000002\) already displayed. In particular \(K>0\).

\subsection{One endpoint controls the whole infinite range}

Writing \(\ell=\log4\), equation \eqref{split-bound} at \(X=q^4\)
becomes
\[
 \frac{S_+(q^4)}L\le
 \frac{\pi\rho}{4u}
 \left(15L+2\ell+2-\frac{\ell^2+2\ell}{L}\right)
 +\frac{2B_q}{u}\left(1-\frac{\ell-1}{L}\right).
\]
Using \(\pi<22/7\) and rigorous bounds for \(\ell\), this is less than
\begin{equation}\label{split-uniform}
 F(q):=\frac{0.4926}{\sqrt q}
 \left(15L+4.773-\frac{4.694}{L}\right)
 +\frac{2}{\sqrt q}\left(\frac{88}{35}+\frac{192}{35\sqrt q}\right)
 \left(1-\frac{0.386}{L}\right).
\end{equation}
Each term of \(F(q)\) decreases for \(L\ge8\). For the first term,
the sign of its derivative with respect to \(L\) is the sign of
\[
 15+\frac{4.694}{L^2}
 -\frac12\left(15L+4.773-\frac{4.694}{L}\right),
\]
which is negative already using \(L\ge8\) in each factor.
For the boundary terms, differentiate
\(e^{-aL}(1-0.386/L)\) with \(a=1/2\) or \(1\); their derivatives
are also negative for \(L\ge8\). Positivity of these factors is
immediate. The other terms in
\(F(q)+4/q+K/L+1/(8L^2)\) decrease as well.

It is therefore enough to evaluate the right side at \(q=4\,000\).
The exact certificate uses only the rational enclosures
\[
 8.294049<\log4000<8.294050,\qquad
 \sqrt{4000}>63.245553.
\]
Substitution gives
\[
 \delta>1.641,\qquad
 F(4000)+\frac4{4000}+\frac{K}{\log4000}
       +\frac1{8(\log4000)^2}<1.634.
\]
The certified gap is greater than \(0.007\).
This contradicts \eqref{necessary} for every \(q\ge4\,000\),
and proves the analytic assertion of Theorem~\ref{main}.

\section{The remaining finite classification}\label{finite}

For \(q>8\), the elementary norm argument in
Section 2 already forces \(q\equiv3\pmod8\) and makes every prime
\(p<q/4\) inert. Hence every remaining candidate is among the 499
integers
\[
 11\le q<4\,000,\qquad q\equiv3\pmod8.
\]
The program \texttt{verify\_finite\_range.py} checks every such integer.
For each odd prime \(p\le109\) with \(4p<q\), it rejects \(q\) if
\(-q\) is zero or a square modulo \(p\). These would respectively
make \(p\) ramified or split, contrary to \eqref{small-inert}.
Including nonfundamental integers in the scan makes the exclusion
stronger; it cannot remove a possible field incorrectly.
The surviving values are exactly
\[
 11,19,43,67,163.
\]
The negative fundamental discriminants with \(q\le8\) give
\(q=3,4,7,8\).

To check class number one for each survivor, use the correspondence
between ideal classes of discriminant \(-q\) and primitive positive
binary quadratic forms \((a,b,c)\) of discriminant
\(b^2-4ac=-q\). Each class has one reduced representative with
\[
 |b|\le a\le c,\qquad
 b\ge0\ \text{if }|b|=a\text{ or }a=c.
\]
The bound \(a\le\sqrt{q/3}\) follows from
\(q=4ac-b^2\ge3a^2\). Enumerating this finite range, retaining
\(\gcd(a,b,c)=1\), gives exactly one reduced form for each of the
nine survivors. Thus the finite program proves the remaining assertion
of Theorem~\ref{main}, conditional only on the analytic argument above
and the standard form-class correspondence.

\section{Reproduction and scope of the checks}

Both verification programs use only the Python standard library.
Run them from the directory containing this paper:
\begin{verbatim}
python3 verify_green_bound.py --sanity
python3 verify_finite_range.py
\end{verbatim}
The first program performs 39 exact arithmetic checks. Logarithms are
bounded by range reduction and the positive series
\[
 \log x=2\sum_{k\ge0}\frac{t^{2k+1}}{2k+1},\qquad
 t=\frac{x-1}{x+1},\quad1\le x\le2.
\]
After 60 terms, its remaining tail is at most
\(2t^{121}/(121(1-t^2))\). Every acceptance decision uses integers
or rational fractions, not floating-point arithmetic.
The optional checks test the valuation identities in 2,880 finite
cases and the derivative row transformations in 60,000 cases,
including \(p=2\) and derivative orders at least \(p\).
The second program performs the complete finite exclusion and
reduced-form calculation.

These checks establish the stated finite calculations and the numerical
gap. They do not certify the general resultant, lattice-counting, or
determinant arguments, nor do they reprove Rosser--Schoenfeld's
theorems or quadratic ideal theory. Those mathematical steps remain
subject to independent review. An exploratory floating-point optimizer
was used to choose parameters; it is not an input to the proof or to
the exact verification programs.

\begin{thebibliography}{9}
\bibitem{green}
B.~Green, \emph{On Landau--Siegel zeros}, version 2, 2026.
Source snapshot identified as \texttt{2610.00311v2},
\href{https://hexagonmath.org/api/v1/works/2610.00311v2/source?path=no-siegel.tex}
{\texttt{no-siegel.tex}}.
The derivative determinant, Frobenius row operation, and resultant
argument used here originate in this note.

\bibitem{rs}
J.~Barkley Rosser and L.~Schoenfeld,
\emph{Approximate formulas for some functions of prime numbers},
Illinois J. Math. \textbf{6} (1962), 64--94.
\href{https://doi.org/10.1215/ijm/1255631807}{doi:10.1215/ijm/1255631807}.
Theorems 6 and 9, printed pp.~70--71.

\bibitem{tao}
T.~Tao, \emph{An experimental AI math challenge: using tenth
discriminant bounds to explore arguments ruling out Siegel zeroes},
20 October 2026.
\href{https://terrytao.wordpress.com/2026/10/20/an-experimental-ai-math-challenge-using-tenth-discriminant-bounds-to-explore-arguments-ruling-out-siegel-zeroes/}
{Challenge and discussion}.

\bibitem{v3}
\emph{A direct split-prime bound for the tenth discriminant},
version 0.3.1, 20 October 2026.
\href{https://github.com/twinheterodox/tenth-discriminant/releases/tag/v0.3.1}
{Archived preceding version}.
The split-prime counting argument in Section 2 is retained from that
version; its Banks determinant is replaced here.
\end{thebibliography}
\end{document}
